% Relatório Técnico do Estudo Dirigido (PPGCC/UFMA, Norma 01/2019, Art. 4º).
% Compilação: latexmk -pdf relatorio_estudo_dirigido.tex
% Mesmo modelo ABNT do relatório de estudo dirigido anterior (packages/abnt-UTP.sty).

\documentclass[
    12pt,
    oneside,
    a4paper,
    chapter=TITLE,
    section=TITLE,
    sumario=abnt-6027-2012,
    english,
    brazil
]{abntex2}
\usepackage[utf8]{inputenc}

\usepackage{packages/abnt-UTP}
\usepackage{amssymb}

\titulo{Estudo Dirigido em Detecção Eficiente de Vulnerabilidades em Código com
Modelos de Linguagem: Quantização e QLoRA sob Avaliação Descontaminada}

\autor{
\normalfont\small
WESLEY DOS SANTOS GATINHO\\
}

\orientador{\normalfont\small
DR. JOÃO OTÁVIO BANDEIRA DINIZ}
\preambulo{Relatório apresentado como requisito para aprovação na disciplina de Estudo Dirigido, do Mestrado em Ciência da Computação da Universidade Federal do Maranhão (UFMA).}

\instituicao{
UNIVERSIDADE FEDERAL DO MARANHÃO \\
Centro de Ciências Exatas e Tecnologia \\
PROGRAMA DE PÓS-GRADUAÇÃO EM CIÊNCIA DA COMPUTAÇÃO \\
}

\local{
\normalfont\small
São Luís - MA}

\data{
\normalfont\small
2026}

\begin{document}

\imprimircapa
\imprimirfolhaderosto{}
\sumario

\textual
\chapter{Introdução}

A detecção de vulnerabilidades em código-fonte pode ser formulada como classificação de funções: dado o corpo de uma função C/C++, um detector indica se ela contém uma vulnerabilidade. Para essa tarefa, os detectores aprendidos passaram de representações do código em grafos \cite{zhou2019devign} a Transformers pré-treinados em código \cite{fu2022linevul} e, mais recentemente, a modelos de linguagem de grande porte (LLMs, do inglês \textit{large language models}) especializados em código, como a família Qwen2.5-Coder \cite{hui2024qwen}.

Dois problemas acompanham essa evolução. O primeiro é de avaliação: duplicação de dados, rótulos imprecisos e distribuição irreal das classes produziram resultados que não se mantêm em condições realistas \cite{chakraborty2020reveal,ding2024primevul}. O segundo é de custo: modelos com bilhões de parâmetros exigem memória e tempo de processamento que limitam seu uso, o que motiva técnicas como a quantização e o ajuste fino eficiente \cite{hu2021lora,dettmers2023qlora}.

Este relatório documenta o Estudo Dirigido realizado no segundo semestre de 2026 no Programa de Pós-Graduação em Ciência da Computação da UFMA, sob supervisão do Prof. Dr. João Otávio Bandeira Diniz. O estudo percorreu os tópicos da ementa aprovada: fundamentos de vulnerabilidades e de modelos de linguagem de código; conjuntos de dados e avaliação descontaminada, com o PrimeVul como estudo de caso; métricas de detecção; ajuste fino eficiente com QLoRA e cabeça de classificação; quantização de modelos; medição de custo computacional; e elaboração de relatório técnico e artigo científico.

Além do estudo teórico, foram executados experimentos com os modelos Qwen2.5-Coder de 3B e 7B parâmetros no PrimeVul. Os resultados deram origem ao artigo ``Limitações das Métricas Globais na Detecção de Vulnerabilidades: Uma Avaliação do Qwen2.5-Coder no PrimeVul'', redigido em português e em inglês para a disciplina de Seminários.

O relatório está organizado da seguinte forma. O Capítulo~2 apresenta os objetivos. O Capítulo~3 reúne a fundamentação teórica dos tópicos estudados. O Capítulo~4 descreve as atividades e a metodologia dos experimentos. O Capítulo~5 apresenta e discute os resultados. O Capítulo~6 traz a conclusão e os próximos passos.

\chapter{Objetivos}

\section{Objetivo Geral}

Estudar a detecção de vulnerabilidades em código-fonte com modelos de linguagem, considerando ao mesmo tempo a qualidade da detecção sob avaliação descontaminada e o custo computacional de técnicas de eficiência, como a quantização e o ajuste fino com QLoRA.

\section{Objetivos Específicos}

Os objetivos específicos do estudo compreenderam:

\begin{alineas}
    \item compreender a formulação da detecção de vulnerabilidades por função e as formas de uso de modelos de linguagem de código nessa tarefa;
    \item estudar os problemas de qualidade dos conjuntos de dados e o protocolo de avaliação do PrimeVul;
    \item estudar e implementar as métricas de detecção: F1, precisão, revocação, taxa de falsos positivos, VD-S e avaliação pareada;
    \item implementar o ajuste fino com QLoRA e cabeça de classificação nos modelos Qwen2.5-Coder de 3B e 7B;
    \item estudar os métodos de quantização NF4, AWQ e GPTQ e avaliar experimentalmente a implantação em NF4;
    \item medir o custo computacional das configurações avaliadas;
    \item registrar os resultados em relatório técnico e em artigo científico.
\end{alineas}

\chapter{Fundamentação Teórica}

\section{Vulnerabilidades e Modelos de Linguagem de Código}

Uma vulnerabilidade é um defeito de software que pode ser explorado para violar a segurança de um sistema, como um acesso à memória fora dos limites ou a falta de validação de uma entrada. Na formulação mais usada pela literatura, a detecção é feita por função: o detector recebe o corpo de uma função e indica se ela é vulnerável. Essa formulação tem uma finalidade delimitada, que é indicar a analistas de segurança quais funções examinar \cite{fu2022linevul}.

Os detectores aprendidos diferem sobretudo na representação do código. O Devign codifica dependências de controle e de dados em um grafo derivado da árvore sintática \cite{zhou2019devign}, enquanto o LineVul substitui o grafo por um Transformer pré-treinado em código \cite{fu2022linevul}. Com os LLMs de código, a tarefa pode ser executada de duas maneiras. Na primeira, uma instrução textual (\textit{prompt}) solicita uma resposta, e um analisador converte o texto em rótulo. Na segunda, uma camada de classificação é acoplada ao modelo e fornece um escore contínuo, que pode ser ordenado e comparado a limiares.

A família Qwen2.5-Coder, usada neste estudo, é oferecida de 0,5B a 32B parâmetros, em variantes Base e Instruct \cite{hui2024qwen}. A especialização em código, porém, não garante a detecção: \citeonline{ullah2024sec} encontraram em oito LLMs respostas não determinísticas, raciocínio infiel e baixa robustez a alterações nas condições de avaliação.

\section{Conjuntos de Dados e Avaliação Descontaminada}

Os conjuntos de dados de vulnerabilidades são, em geral, construídos a partir de commits de correção: a versão anterior ao commit é rotulada como vulnerável e a posterior, como corrigida. Esse procedimento introduz três problemas. O primeiro é o ruído nos rótulos, pois nem toda função alterada em um commit de correção é vulnerável \cite{chen2023diversevul}. O segundo é a duplicação, que permite a uma mesma função aparecer no treino e no teste. O terceiro é a distribuição irreal das classes, com proporção de funções vulneráveis muito acima da encontrada em projetos reais \cite{chakraborty2020reveal}.

Uma avaliação é dita descontaminada quando o conjunto de teste não contém informação vista no treino. O PrimeVul foi construído com esse objetivo, por meio de rotulação mais rigorosa, deduplicação e divisão cronológica entre treino, validação e teste; nesse conjunto, um modelo de 7B que alcançava F1 elevado no BigVul obteve F1 muito baixo \cite{ding2024primevul}.

O PrimeVul também fornece pares formados por uma função vulnerável e sua versão corrigida. Como as duas versões compartilham projeto e contexto, um detector que as distingue dificilmente está usando uma característica alheia à vulnerabilidade. Esse desenho responde a um achado recorrente: modelos treinados em conjuntos típicos não distinguem funções vulneráveis de suas correções \cite{risse2024limits}, e LLMs ajustados alcançam escores altos contra código não relacionado, mas falham diante das versões corrigidas \cite{huang2026trap}.

\section{Métricas de Detecção}

Sejam VP, FP, VN e FN os verdadeiros positivos, falsos positivos, verdadeiros negativos e falsos negativos, com a classe vulnerável como positiva. As métricas globais usadas são:

\begin{equation}
\text{Precisão} = \frac{VP}{VP + FP}, \qquad
\text{Revocação} = \frac{VP}{VP + FN}
\end{equation}

\begin{equation}
F1 = \frac{2 \cdot \text{Precisão} \cdot \text{Revocação}}{\text{Precisão} + \text{Revocação}}
\end{equation}

\begin{equation}
FPR = \frac{FP}{FP + VN}, \qquad FNR = \frac{FN}{FN + VP}
\end{equation}

Como a classe vulnerável é rara, a acurácia é pouco informativa. A área sob a curva ROC (AUC) também exige cuidado, pois em dados desbalanceados ela pode ser alta mesmo com precisão baixa \cite{saito2015metrics}.

O \textit{Vulnerability Detection Score} (VD-S), proposto com o PrimeVul, mede a taxa de falsos negativos quando a taxa de falsos positivos é limitada a um valor baixo, como 0,5\% \cite{ding2024primevul}. A métrica representa o uso prático, em que um analista só tolera poucos alarmes falsos.

Na avaliação pareada, cada par recebe uma de quatro categorias, segundo as decisões sobre as duas versões: P-C, ambas corretas; P-V, ambas classificadas como vulneráveis; P-B, ambas classificadas como não vulneráveis; e P-R, decisões invertidas. Quando o detector produz um escore, é possível também medir a ordenação pareada, isto é, a frequência com que a versão vulnerável recebe escore maior que a corrigida, com meio crédito para empates:

\begin{equation}
A_{\text{par}} = \frac{1}{N}\sum_{j=1}^{N}\left[\mathbf{1}(s_{v,j} > s_{c,j}) + \tfrac{1}{2}\,\mathbf{1}(s_{v,j} = s_{c,j})\right]
\end{equation}

em que $s_{v,j}$ e $s_{c,j}$ são os escores da versão vulnerável e da corrigida no par $j$. Um escore sem relação com as versões tende a 50\%.

\section{Ajuste Fino Eficiente com QLoRA}

Na adaptação de baixo posto (LoRA), a atualização de uma matriz pré-treinada $W_0 \in \mathbb{R}^{d \times k}$ é representada pelo produto de duas matrizes de baixo posto, de modo que o peso adaptado é $W_0 + BA$, com $B \in \mathbb{R}^{d \times r}$, $A \in \mathbb{R}^{r \times k}$ e $r \ll \min(d,k)$. Durante o treino, apenas $A$ e $B$ são atualizadas, o que reduz o número de parâmetros treináveis a uma pequena fração do modelo \cite{hu2021lora}. Como a atualização é linear, o adaptador pode ser mesclado aos pesos originais após o treino.

O QLoRA combina a LoRA com uma base quantizada em quatro bits: os pesos originais ficam congelados no formato NF4, e os gradientes atravessam a base quantizada para atualizar apenas os adaptadores, mantidos em precisão maior \cite{dettmers2023qlora}. Com isso, o ajuste de um modelo de 7B cabe em uma única GPU.

Para classificação, acrescenta-se ao modelo uma cabeça de duas saídas sobre a representação do último token da sequência. O escore de vulnerabilidade é a diferença entre os dois logits, e a cabeça é treinada junto com os adaptadores.

\section{Quantização de Modelos}

A quantização representa os pesos de um modelo com menos bits, reduzindo memória e, conforme o suporte de hardware, o tempo de inferência. Na quantização pós-treinamento (PTQ), o modelo já treinado é convertido sem novo treino, eventualmente com um pequeno conjunto de calibração \cite{gholami2021survey}. Os três métodos da ementa são de quatro bits e diferem na forma de escolher os valores quantizados:

\begin{alineas}
    \item \textbf{NF4} (\textit{4-bit NormalFloat}): parte da observação de que os pesos pré-treinados seguem aproximadamente uma distribuição normal e define os 16 níveis pelos quantis dessa distribuição. Não usa dados de calibração e pode ser aplicado no momento da carga do modelo. A dupla quantização reduz ainda o custo das constantes de escala \cite{dettmers2023qlora}.
    \item \textbf{GPTQ}: quantiza os pesos camada por camada, usando informação aproximada de segunda ordem (a Hessiana do erro de reconstrução da camada). A cada peso quantizado, os pesos restantes são ajustados para compensar o erro introduzido. Exige dados de calibração \cite{frantar2023gptq}.
    \item \textbf{AWQ}: parte da observação de que os pesos não são igualmente importantes e identifica os mais relevantes pela magnitude das ativações. Esses pesos são protegidos por um fator de escala por canal, buscado em dados de calibração, sem retropropagação \cite{lin2024awq}.
\end{alineas}

Há, portanto, duas formas de obter um classificador de quatro bits: ajustar o modelo sobre a base NF4 e mantê-la na implantação, ou ajustar, mesclar o adaptador e só então quantizar com AWQ ou GPTQ. A comparação entre essas formas depende de medir se a discriminação se mantém em cada uma.

\section{Medição de Custo Computacional}

O custo de uma configuração foi estudado em três dimensões. A memória pode ser medida como a memória de GPU alocada após a carga do modelo ou como o pico durante a inferência; as duas medidas diferem da memória reservada e da capacidade da placa. A latência pode ser medida por requisição individual ou como tempo amortizado por função em lotes; os dois valores não são intercambiáveis. A energia é obtida pela integração da potência da GPU ao longo da execução e depende do equipamento medido.

Em todas as dimensões, a medida só é comparável quando feita no mesmo tipo de GPU, com as mesmas versões de bibliotecas e o mesmo tamanho de lote. O custo financeiro é derivado do tempo, multiplicado pelo preço da hora de GPU.

\chapter{Metodologia}

\section{Atividades do Estudo}

A Tabela~\ref{tab:ementa} relaciona cada tópico da ementa à atividade realizada. Os tópicos foram cobertos por leitura dos trabalhos de referência e, quando indicado, por implementação e experimento.

\begin{table}[H]
    \centering
    \caption{Tópicos da ementa e atividades realizadas.}
    \label{tab:ementa}
    \footnotesize
    \begin{tabular}{|p{4.6cm}|p{9.4cm}|}
        \hline
        \textbf{Tópico da ementa} & \textbf{Atividade realizada} \\
        \hline
        Fundamentos de vulnerabilidades e de modelos de linguagem de código & Estudo da formulação por função e dos detectores Devign, LineVul e LLMs de código; uso do Qwen2.5-Coder por instrução e por cabeça de classificação. \\
        \hline
        Datasets e avaliação descontaminada (PrimeVul) & Estudo dos problemas de rotulação, duplicação e divisão dos dados; uso da distribuição original do PrimeVul, com conferência dos totais e dos pares. \\
        \hline
        Métricas de detecção (F1, VD-S e avaliação pareada) & Implementação de F1, precisão, revocação, FPR, AUC, FNR com FPR limitada, categorias pareadas e ordenação pareada. \\
        \hline
        Ajuste fino eficiente com QLoRA e cabeça de classificação & Treino de seis classificadores (três receitas no 3B e três sementes no 7B). \\
        \hline
        Quantização de modelos (NF4, AWQ e GPTQ) & Estudo teórico dos três métodos; avaliação experimental da implantação em NF4, comparada à implantação em BF16. A comparação experimental com AWQ e GPTQ não foi executada neste estudo. \\
        \hline
        Medição de custo computacional (memória, latência e energia) & Medição de memória, tempo por função e custo estimado. A energia foi estudada, mas não foi medida. \\
        \hline
        Elaboração de relatório técnico e artigo científico & Este relatório e o artigo da disciplina de Seminários, em português e em inglês. \\
        \hline
    \end{tabular}
\end{table}

\section{Conjunto de Dados}

O estudo utiliza a distribuição original do PrimeVul, composta por funções C/C++ de projetos de código aberto, com rótulos derivados de commits de correção \cite{ding2024primevul}. A Tabela~\ref{tab:particoes} apresenta as partições e o número de pares vulnerável/corrigida em cada uma.

\begin{table}[H]
    \centering
    \caption{Partições do PrimeVul utilizadas.}
    \label{tab:particoes}
    \begin{tabular}{lrrr}
        \toprule
        \textbf{Partição} & \textbf{Funções} & \textbf{Vulneráveis} & \textbf{Pares} \\
        \midrule
        Treino & 184.427 & 5.574 & 4.354 \\
        Validação & 25.430 & 699 & 562 \\
        Teste & 25.911 & 695 & 564 \\
        \bottomrule
    \end{tabular}
\end{table}

A classe vulnerável corresponde a 2,68\% do teste global. Cada par reúne uma função rotulada como vulnerável e outra rotulada como sua versão corrigida, no mesmo projeto e commit. Os 564 pares de teste pertencem a 550 grupos projeto/commit.

\section{Configurações Avaliadas}

São avaliados três tipos de configuração:

\begin{alineas}
    \item \textbf{Classificação por instrução.} As variantes Qwen2.5-Coder-3B-Instruct e 7B-Instruct, carregadas em 16 bits, recebem uma instrução sem exemplos que pede uma única palavra, YES ou NO, para cada função. A decodificação é determinística, e um analisador aceita a primeira palavra de decisão reconhecida. O contexto de entrada é de até 8.192 tokens.
    \item \textbf{Heurística de comprimento.} Atribui a cada função um escore igual ao número de tokens visíveis após o corte em 2.048 tokens. Não usa nenhuma informação sobre a vulnerabilidade e indica quanto uma característica superficial correlacionada ao rótulo alcança em cada métrica.
    \item \textbf{Classificadores com QLoRA.} Partem dos modelos Base Qwen2.5-Coder-3B e 7B, com cabeça de duas saídas e contexto de 2.048 tokens. No ajuste, a base é carregada em NF4 com dupla quantização e computação em BF16. A LoRA é aplicada às projeções de atenção e do bloco MLP, com posto $r = 16$, $\alpha = 32$ e \textit{dropout} de 0,05. O otimizador é o AdamW, com taxa de aprendizado inicial de $2 \times 10^{-4}$, aquecimento em 3\% dos passos e decaimento linear. O lote efetivo tem 16 funções, e a perda é a entropia cruzada.
\end{alineas}

No 3B, são exploradas três receitas de amostragem do treino. A primeira mantém todas as funções vulneráveis e sorteia igual número de não vulneráveis (1:1). A segunda usa o treino completo, na proporção original. A terceira forma o conjunto 1:1 escolhendo, para cada função vulnerável, a não vulnerável de comprimento mais próximo (1:1 com controle de tamanho). O 7B usa somente essa última receita, com três sementes de aleatoriedade. A Tabela~\ref{tab:treinos} resume os treinos.

\begin{table}[H]
    \centering
    \caption{Treinos realizados: funções vulneráveis e não vulneráveis usadas, épocas executadas, época selecionada e horas de treino.}
    \label{tab:treinos}
    \footnotesize
    \begin{tabular}{lrrrrr}
        \toprule
        \textbf{Configuração} & \textbf{Vuln.} & \textbf{Não vuln.} & \textbf{Épocas} & \textbf{Selecionada} & \textbf{Treino (h)} \\
        \midrule
        3B 1:1 & 5.574 & 5.574 & 3 & 1 & 3,37 \\
        3B proporção original & 5.574 & 178.853 & 1 & 1 & 10,16 \\
        3B 1:1 + tamanho & 5.574 & 5.574 & 3 & 3 & 4,98 \\
        7B 1:1 + tamanho, semente 1 & 5.574 & 5.574 & 3 & 3 & 7,89 \\
        7B 1:1 + tamanho, semente 2 & 5.574 & 5.574 & 3 & 2 & 7,89 \\
        7B 1:1 + tamanho, semente 3 & 5.574 & 5.574 & 3 & 3 & 7,91 \\
        \bottomrule
    \end{tabular}
\end{table}

\section{Seleção de Época, Limiares e Implantação}

A época de cada classificador é escolhida pela ordenação pareada nos 562 pares de validação. São calibrados dois limiares na validação global: o primeiro maximiza F1; o segundo é o menor limiar que mantém a FPR em no máximo 0,5\% e serve para medir a FNR no teste, em aproximação ao VD-S. Ambos são aplicados ao teste sem recalibração.

Cada classificador 7B selecionado é avaliado em duas implantações. Na implantação NF4, a base permanece quantizada e o adaptador separado. Na implantação BF16, a base é carregada em BF16 e o adaptador é mesclado aos pesos. Como as duas diferem ao mesmo tempo na representação da base e na incorporação do adaptador, elas são tratadas como configurações completas de implantação, e não como um experimento que isola a quantização.

\section{Medição de Recursos e Análise Estatística}

Os recursos são registrados em uma GPU NVIDIA A40, com lotes de até 32.768 tokens. A memória corresponde à memória CUDA alocada após a carga do modelo, e o pico é observado durante a pontuação. O tempo por função é o tempo de pontuação do teste dividido pelo número de funções, ou seja, um tempo amortizado em lotes. O custo estimado multiplica esse tempo pelo preço de US\$ 0,50 por hora de GPU.

No 7B, a variação entre treinamentos é descrita pela média e pelo desvio padrão das três sementes. A incerteza amostral da ordenação pareada é estimada por \textit{bootstrap} percentil de 95\% com 10.000 réplicas, reamostrando os 550 grupos projeto/commit com reposição \cite{field2007bootstrap}.

\section{Ambiente e Reprodutibilidade}

Os experimentos são executados em instâncias alugadas com uma GPU NVIDIA A40 de 48 GB, com PyTorch e Transformers em versões travadas em arquivo de requisitos. Cada execução registra em planilha as versões das bibliotecas, a revisão do modelo e a impressão digital dos dados. O código, os registros de execução e as planilhas de resultados estão versionados no repositório do projeto.

\chapter{Resultados e Discussão}

\section{Classificação por Instrução}

A Tabela~\ref{tab:instrucao} apresenta o desempenho dos modelos instruídos por \textit{prompt}, sem ajuste. O F1 foi baixo nas duas escalas, e o modelo maior usou mais memória e mais tempo sem ganho de qualidade: sua revocação um pouco maior veio acompanhada de uma taxa de falsos positivos três vezes maior.

\begin{table}[H]
    \centering
    \caption{Classificação por instrução no teste global (25.911 funções) e nos 564 pares. Memória de GPU e tempo por função medidos no teste global.}
    \label{tab:instrucao}
    \footnotesize
    \begin{tabular}{lrrrrrrr}
        \toprule
        \textbf{Modelo} & \textbf{F1} & \textbf{Precisão} & \textbf{Revocação} & \textbf{FPR (\%)} & \textbf{P-C} & \textbf{Mem. (GB)} & \textbf{Tempo (s)} \\
        \midrule
        3B-Instruct & 0,119 & 0,102 & 0,144 & 3,49 & 16 & 7,2 & 0,071 \\
        7B-Instruct & 0,062 & 0,038 & 0,160 & 11,00 & 14 & 16,9 & 0,112 \\
        \bottomrule
    \end{tabular}
\end{table}

Nos pares, a categoria P-B reuniu 463 e 454 dos 564 casos. A taxa de respostas YES foi praticamente a mesma na versão vulnerável e na corrigida: 14,7\% e 15,1\% no 3B, e 16,3\% e 17,0\% no 7B. A resposta dos modelos, portanto, não acompanhou a correção da falha. Esses resultados descrevem a instrução utilizada, e não todo o espaço de instruções possíveis.

\section{Desempenho Global e Ordenação Pareada}

A Tabela~\ref{tab:global} reúne, para as configurações com escore, o desempenho no teste global e a ordenação nos 564 pares.

\begin{table}[H]
    \centering
    \caption{Teste global e teste pareado. F1 e P-C usam o limiar de maior F1 na validação. FNR* usa o limiar calibrado para FPR $\leq$ 0,5\% na validação. s1--s3 indicam as sementes do 7B.}
    \label{tab:global}
    \footnotesize
    \begin{tabular}{lrrrrr}
        \toprule
        \textbf{Configuração} & \textbf{F1} & \textbf{AUC} & \textbf{FNR* (\%)} & \textbf{P-C} & \textbf{Ordenação pareada (\%)} \\
        \midrule
        Heurística de comprimento & 0,200 & 0,810 & 100,00 & 2 & 21,10 \\
        3B 1:1 & 0,195 & 0,795 & 94,10 & 13 & 42,38 \\
        3B proporção original & 0,262 & 0,851 & 88,06 & 6 & 37,50 \\
        3B 1:1 + tamanho & 0,121 & 0,746 & 96,26 & 6 & 49,20 \\
        7B NF4 s1 & 0,114 & 0,738 & 97,99 & 7 & 55,23 \\
        7B BF16 s1 & 0,119 & 0,740 & 98,13 & 6 & 53,28 \\
        7B NF4 s2 & 0,142 & 0,726 & 98,71 & 5 & 45,04 \\
        7B BF16 s2 & 0,148 & 0,728 & 98,71 & 7 & 44,68 \\
        7B NF4 s3 & 0,127 & 0,765 & 98,27 & 6 & 51,60 \\
        7B BF16 s3 & 0,124 & 0,765 & 98,27 & 8 & 51,95 \\
        \bottomrule
    \end{tabular}
\end{table}

O contraste mais nítido envolve a heurística de comprimento. Usando apenas o número de tokens, ela alcançou AUC global de 0,810 e F1 de 0,200, valores próximos aos do classificador 3B 1:1 e acima dos dois modelos instruídos. Nos pares, porém, sua ordenação foi de apenas 21,10\%. Assim, a característica que separa as classes no teste global não separa, e tende a inverter, as duas versões de uma mesma correção.

Os classificadores repetem esse contraste em menor grau. As duas configurações de maior F1 e AUC, a de proporção original e a 1:1, tiveram ordenações pareadas de 37,50\% e 42,38\%. As configurações com controle de tamanho ficaram entre 44,68\% e 55,23\%, com F1 entre 0,114 e 0,148. Esse padrão é compatível com a hipótese de que parte do desempenho global dependa do comprimento da função, embora o desenho não permita identificar essa causa.

Com a FPR limitada a 0,5\% na validação, a FNR no teste foi de 88,06\% na melhor configuração e ficou entre 97,99\% e 98,71\% no 7B. Em um cenário de poucos alarmes falsos, portanto, todas as configurações perderam a maior parte das funções vulneráveis.

\section{Variação entre Sementes no 7B}

Nas três sementes do 7B, a ordenação pareada média foi de 50,62\% na implantação NF4 e de 49,97\% na BF16, com desvios padrão entre sementes de 5,17 e 4,63 pontos percentuais. As sementes isoladas apontam em direções opostas: a semente 1 ficou acima de 50\% nas duas implantações, e a semente 2, abaixo. Uma única semente poderia, assim, sugerir conclusões contrárias. A diferença NF4 $-$ BF16 foi de 0,65 ponto percentual, com intervalo de 95\% de $-1{,}22$ a $2{,}54$, que inclui valores nas duas direções e não distingue as implantações.

\section{Recursos das Implantações}

A Tabela~\ref{tab:recursos} apresenta os recursos medidos nas duas implantações do 7B.

\begin{table}[H]
    \centering
    \caption{Recursos do 7B na NVIDIA A40: memória alocada após a carga, pico durante a pontuação, tempo amortizado por função e custo estimado da pontuação do teste. Faixas entre as três sementes.}
    \label{tab:recursos}
    \footnotesize
    \begin{tabular}{lrrrr}
        \toprule
        \textbf{Implantação} & \textbf{Carga (GB)} & \textbf{Pico (GB)} & \textbf{Tempo (ms)} & \textbf{Custo (US\$)} \\
        \midrule
        NF4, adaptador separado & 4,63 & 12,49 & 59,6--60,0 & 0,2145--0,2160 \\
        BF16, adaptador mesclado & 14,15 & 20,76 & 46,8--47,1 & 0,1684--0,1696 \\
        \bottomrule
    \end{tabular}
\end{table}

A implantação NF4 alocou cerca de 67\% menos memória após a carga e cerca de 40\% menos no pico. Em contrapartida, seu tempo por função foi maior, e o custo estimado acompanhou essa diferença. Nesse ambiente, a implantação NF4 usou menos memória e mais tempo, sem diferença distinguível na ordenação pareada.

\section{Reprodutibilidade}

Durante a preparação dos experimentos, a mesma rodada de classificação por instrução foi repetida em condições controladas, em uma versão anterior e menor do conjunto de pares (870 funções). Em uma instância nova com a mesma GPU e as mesmas versões de bibliotecas, as 870 respostas foram idênticas. Em uma GPU de outro modelo, 9 das 870 respostas mudaram. Por esse motivo, todas as comparações deste estudo foram feitas no mesmo tipo de GPU e com as versões travadas.

\section{Discussão}

O principal resultado é a dissociação entre desempenho global e ordenação pareada. Uma heurística sem qualquer informação sobre vulnerabilidades obteve AUC global de 0,810, e o classificador com maior F1 teve ordenação pareada abaixo de 50\%. Esse quadro é coerente com a observação de \citeonline{chakraborty2020reveal} de que detectores profundos frequentemente não aprendem características da vulnerabilidade e com o achado de \citeonline{risse2024limits} sobre funções vulneráveis e suas correções.

Quanto ao custo, o estudo mostra que a escolha entre implantações precisa ser lida junto à qualidade: a implantação em quatro bits reduziu a memória, mas não o tempo, e nenhuma das duas distinguiu as versões dos pares.

Os resultados têm limitações. O conjunto de teste foi consultado durante o desenvolvimento das receitas de treino, de modo que a análise é exploratória, e não confirmatória. Cada configuração 3B tem uma única semente, e o 7B tem três, quantidade pequena diante da dispersão observada. Além disso, os resultados se referem a uma família de modelos, um conjunto de dados e uma GPU.

\chapter{Conclusão}

O Estudo Dirigido percorreu os tópicos da ementa, da formulação do problema à medição de custo, e produziu um conjunto de experimentos reprodutíveis com os modelos Qwen2.5-Coder de 3B e 7B no PrimeVul. O domínio adquirido abrange a avaliação descontaminada de detectores, a implementação das métricas globais e pareadas, o ajuste fino com QLoRA e cabeça de classificação, os fundamentos dos métodos de quantização de quatro bits e a medição de memória e tempo.

Os resultados sustentam três observações para as configurações executadas. Primeiro, desempenho global e discriminação pareada não andaram juntos: a heurística de comprimento e o classificador de maior F1 tiveram ordenação pareada abaixo de 50\%. Segundo, a variação entre as três sementes do 7B foi grande em relação às diferenças entre implantações. Terceiro, a implantação NF4 usou menos memória e mais tempo que a BF16, sem diferença distinguível na ordenação.

Dois itens da ementa foram cobertos apenas no plano teórico e ficam como continuidade: a comparação experimental entre NF4, AWQ e GPTQ sobre o modelo ajustado e a medição de energia. Também é necessária uma avaliação confirmatória, com receitas, critérios e limiares definidos antes do acesso ao teste e mais sementes por configuração. Essas etapas integram o projeto de mestrado ao qual este estudo está vinculado.

Como produtos do estudo, foram elaborados este relatório técnico e o artigo científico apresentado à disciplina de Seminários, em versões em português e em inglês.


\postextual
\bibliography{referencias/referencias}

\end{document}
